{\Large\textbf{1\quad Light-emitting diode}}

Light Emitting Diodes (LEDs) are the most efficient light sources. During
recent years, the cheap, powerful and reliable LED light bulbs appeared on the
market. This is making a true revolution worldwide, as other electrical
lighting tools (e.g. incandescent, halogen and fluorescent) in homes and
offices are being replaced with LED lighting.

In this experiment, we will analyse thermal and electrical properties of the
light emitting diodes.

You do not need to estimate any uncertainties, but the accuracy of your methods
and results is important and will be graded. Always draw the measurement setup
you use! When appropriate, use graphs for determining required quantities.

\textit{Equipment: 2 identical circuit boards with a LED, resistor and a
temperature sensor on them; 2 transparent bottles, 2 airtight caps, 2 tubes,
water, syringe, 3 multimeters (the multimeter labelled as ``voltage-only'' is
to be used only for voltage measurements), 2 potentiometers, power supply,
cables, stand, stripes of graph paper with millimeter scale.}

\begin{center}\includegraphics[width=0.50\linewidth]{figures/EuPhO_2017_Q4_fig1.png}\end{center}

The schematic and connections of the circuit board (left) and the
potentiometer (right) are in the figure. The resistor $R_1$ may be used for
heating the board near the LED. The temperature sensor (thermistor) $R_T$ is
another resistor whose resistance depends strongly on the absolute temperature
$T$:
\[
T = 2.254\,\mathrm{K}\left(\ln\frac{R_T}{1\,\mathrm{k\Omega}}\right)^2
- 32.46\,\mathrm{K}\ln\frac{R_T}{1\,\mathrm{k\Omega}} + 361.09\,\mathrm{K}.
\]

\textbf{Warning! Apply voltage to the LED only with the polarity shown!} The
red lead of the power supply is ``+'' (and should be connected to the red
connector of the LED) and the black lead is ``$-$''.

The multimeter has a mode marked
``$\scriptstyle-\!\!\blacktriangleright\!\!|\!\!-$'' that acts as a source of
an approximately constant small current, when a diode is connected between the
``mAV$\Omega$'' (supplies ``+'') and ``COM'' (``$-$'') terminals. In this mode
the multimeter shows the voltage across the diode in volts, while supplying
about $0.33\,\mathrm{mA}$ (you may assume the current stays fixed).

In a simplified theory, the diode current $I_d$, voltage drop $V$ at the
junction inside the LED where light is emitted, and absolute temperature $T$ of
the junction obey
\[
I_d = Ae^{-V_{G0}/(nV_T)}\left(e^{V/(nV_T)} - 1\right) \quad \text{with } V_T = \frac{k}{q}T,
\]
Boltzmann constant $k = 1.381 \times 10^{-23}\,\mathrm{J/K}$ and elementary
charge $q = 1.602 \times 10^{-19}\,\mathrm{C}$. The variable $V_T$ is called
the thermal voltage. The parameters $V_{G0}$ and $n$ depend on the materials of
the LED; the parameter $A$ depends also on the construction of the LED. The
parameter $n$ is called the ideality factor and usually $1 < n < 2$. The
parameter $V_{G0}$ is called the nominal bandgap voltage of the semiconductor
material.

The voltage across a physical diode $V' = V + I_dR_s$ has also contribution
from a parasitic series resistance $R_s$ which is of the order of
$1\,\Omega$. \textbf{Hint:} estimate the magnitudes in the expression above and
simplify your calculations accordingly!

1. Measure and plot the voltage--temperature graph of the LED at a constant
current (your current should be small enough so that the voltage drop on $R_s$
can be neglected).

Find $V_{G0}$.

Find the parameters $n$ and $A$ by making additional measurements and a
suitable plot.

At larger currents, the series resistance $R_s$ becomes noticeable. Measure
this $R_s$.

2. Define the efficiency of the LED as the ratio between the power radiated as
light and the consumed electrical power. Measure a value of the efficiency
$\eta$ of the LED without using the temperature sensor.

3. The LED also behaves as a solar cell (or a photodiode). The photocurrent
$I_p$ generated by light does not depend on the voltage and is proportional to
the light intensity; it is subtracted from the diode current
($I = I_d - I_p$). The photocurrent from the ambient light was low enough not
to affect previous measurements.

Place two LEDs directly opposite to each other at $d = 3.0\,\mathrm{cm}$
distance, and supply one of them with $I_1 = 0.50\,\mathrm{A}$. Determine the
maximum electrical power $P_{\mathrm{max}}$ that can be harvested from the LED
in this lighting setup at room temperature.

Determine the corresponding photoefficiency $\eta_p$ --- electrical power
output divided by the power of the light absorbed by the active area of the
LED. This area is $S = 1.56\,\mathrm{mm}^2$. Assume that the LED radiates
uniformly into $\alpha = 33\%$ of the sphere.
